It remains to prove that ${}^tu(R'^*)\subset R^*$, which follows from the fact that the elements of $M^*$ (resp. $M'^*$) are determined by duality with $R$ or $\Delta$ (resp. $R'$ or $\Delta'$), by 1.2.5.
\begin{corollary}{7.3.3}\label{XXI.7.3.3}
A reduced simply connected or adjoint root datum is determined up to isomorphism by its Cartan matrix.
\end{corollary}
\begin{corollary}{7.3.4}\label{XXI.7.3.4}
Let $\mathcal R$ be a reduced simply connected (resp. adjoint) root datum, and $\Delta$ a system of simple roots. The group $E_\Delta(\mathcal R)$ is identified with the group of automorphisms of the set $\Delta$ which leave the Cartan matrix invariant.
\end{corollary}
\begin{remark}{7.3.5}\label{XXI.7.3.5}
The question of the existence of a root datum corresponding to a given Cartan matrix satisfying (i), (ii) and (iv) (for example) is not easily solved directly, without using the classification.
\end{remark}

\sgasection{7.4}{Dynkin diagrams}
\begin{definition}{7.4.1}\label{XXI.7.4.1}
A \emph{Dynkin diagram structure} (the word ``scheme'' has been banished for obvious reasons) on a finite set $\Delta$ is the datum of a set of pairs of distinct elements of $\Delta$, called \emph{linked pairs}, and of a map from $\Delta$ to the set $\{1,2,3\}$. The notion of isomorphism of such structures is obvious.
\end{definition}
\begin{definition}{7.4.2}\label{XXI.7.4.2}
Let $\mathcal R$ be a root datum and $\Delta$ a system of simple roots. One calls the \emph{Dynkin diagram} of $\mathcal R$ relative to $\Delta$ the set $\Delta$, two simple roots being linked if and only if they are not orthogonal, each root being assigned its length.
\end{definition}
\begin{proposition}{7.4.3}\label{XXI.7.4.3}
Dynkin diagram and Cartan matrix determine one another bijectively.
\end{proposition}
Indeed, the equivalence
\[
\alpha\text{ and }\beta\text{ are not linked}\Longleftrightarrow a_{\alpha,\beta}=0,
\]
and the relation
\[
\operatorname{long}(\alpha)a_{\alpha,\beta}=\operatorname{long}(\beta)a_{\beta,\alpha}
\]
(with $\inf\operatorname{long}(\alpha)=1$ in each connected component of the diagram) determine the $a_{\alpha,\beta}$ in terms of the links and lengths, and conversely (the details of the verification are left to the reader).
\begin{corollary}{7.4.4}\label{XXI.7.4.4}
A reduced simply connected or adjoint root datum is determined by its Dynkin diagram.
\end{corollary}
\begin{corollary}{7.4.5}\label{XXI.7.4.5}
Let $\mathcal R$ be a reduced simply connected or adjoint root datum and $\Delta$ a system of simple roots. The group $E_\Delta(\mathcal R)$ is identified with the automorphism group of the Dynkin diagram of $\mathcal R$ relative to $\Delta$, that is, with the group of permutations of $\Delta$ preserving lengths and links.
\end{corollary}
\begin{remark}{7.4.6}\label{XXI.7.4.6}
By the usual method,{\renewcommand{\thefootnote}{($*$)}\footnote{Cf. Bourbaki, \textit{Groupes et alg\`ebres de Lie}, chap.~VI \No~4.2, or S\'eminaire Sophus Lie.}}{\renewcommand{\thefootnote}{(30)}\nde{For a slightly different proof, see also [De80].}} one classifies the various connected Dynkin diagrams and shows that each actually corresponds to a reduced simply connected irreducible root datum. One finds the well-known types:
\begin{center}
\begin{tabular}{ll}
$\xymatrix@C=0.55em{\overset{1}{\bullet}\ar@{-}[r]&\overset{1}{\bullet}\ar@{-}[r]&\overset{1}{\bullet}&\cdots&\overset{1}{\bullet}\ar@{-}[r]&\overset{1}{\bullet}\ar@{-}[r]&\overset{1}{\bullet}}$&$A_n,\quad n\geq1.$\\[1em]
$\xymatrix@C=0.55em{\overset{2}{\bullet}\ar@{-}[r]&\overset{2}{\bullet}\ar@{-}[r]&\overset{2}{\bullet}&\cdots&\overset{2}{\bullet}\ar@{-}[r]&\overset{2}{\bullet}\ar@{-}[r]&\overset{1}{\bullet}}$&$B_n,\quad n\geq2.$\\[1em]
$\xymatrix@C=0.55em{\overset{1}{\bullet}\ar@{-}[r]&\overset{1}{\bullet}\ar@{-}[r]&\overset{1}{\bullet}&\cdots&\overset{1}{\bullet}\ar@{-}[r]&\overset{1}{\bullet}\ar@{-}[r]&\overset{2}{\bullet}}$&$C_n,\quad n\geq3.$\\[1em]
$\xymatrix@C=0.55em@R=0.4em{&&&&&\overset{1}{\bullet}\\
\overset{1}{\bullet}\ar@{-}[r]&\overset{1}{\bullet}\ar@{-}[r]&\overset{1}{\bullet}&\cdots&\overset{1}{\bullet}\ar@{-}[ur]\ar@{-}[dr]&\\
&&&&&\overset{1}{\bullet}}$&$D_n,\quad n\geq4.$\\[1em]
$\xymatrix@C=0.55em@R=0.7em{\overset{1}{\bullet}\ar@{-}[r]&\overset{1}{\bullet}&\cdots&\overset{1}{\bullet}\ar@{-}[r]\ar@{-}[d]&\overset{1}{\bullet}\ar@{-}[r]&\overset{1}{\bullet}\\
&&&\overset{1}{\bullet}&&}$&$E_n,\quad n=6,7,8.$\\[1em]
$\xymatrix@C=0.55em{\overset{1}{\bullet}\ar@{-}[r]&\overset{1}{\bullet}\ar@{-}[r]&\overset{2}{\bullet}\ar@{-}[r]&\overset{2}{\bullet}}$&$F_4.$\\[1em]
$\xymatrix@C=0.55em{\overset{1}{\bullet}\ar@{-}[r]&\overset{3}{\bullet}}$&$G_2.$
\end{tabular}
\end{center}
By 7.4.5, one immediately finds the corresponding group $E_\Delta(\mathcal R)$; one has:
\[
\begin{aligned}
E_\Delta(\mathcal R)&=\{e\}&&\text{for }A_1,B_n,C_n,E_7,E_8,F_4,G_2.\\
E_\Delta(\mathcal R)&=\ZZ/2\ZZ&&\text{for }A_n\ (n\geq2),D_n\ (n\geq5),E_6.\\
E_\Delta(\mathcal R)&=\mathfrak S_3&&\text{for }D_4.
\end{aligned}
\]
\end{remark}

\sgasection{7.5}{Supplements on $p$-morphisms}
Let $f:\mathcal R\to\mathcal R'$ be a $p$-morphism (cf. 6.8). It is clear from the definitions that the associated bijection $u:R\isomto R'$ makes the systems of simple roots, systems of positive roots, irreducible components (etc.) of $R$ and $R'$ correspond. Suppose therefore, for simplicity, that $R$ and $R'$ are \emph{irreducible}.
\begin{lemma}{7.5.1}\label{XXI.7.5.1}
If $R$ and $R'$ are irreducible, there exists $k\in\QQ$ such that for every $\alpha\in R$
\[
k\operatorname{long}(u(\alpha))=q(\alpha)^2\operatorname{long}(\alpha).
\]
\end{lemma}
Indeed, one has $\operatorname{long}(\alpha)(\alpha^*,\beta)=\operatorname{long}(\beta)(\beta^*,\alpha)$ and, similarly,
\[
\operatorname{long}(u(\alpha))(u(\alpha)^*,u(\beta))
=\operatorname{long}(u(\beta))(u(\beta)^*,u(\alpha)).
\]
One deduces at once that for nonorthogonal $\alpha$ and $\beta$, one has
\[
\frac{q(\alpha)^2\operatorname{long}(\alpha)}{\operatorname{long}(u(\alpha))}
=\frac{q(\beta)^2\operatorname{long}(\beta)}{\operatorname{long}(u(\beta))},
\]
and one then concludes by 7.1.3 (iv).
\begin{remark}{7.5.2}\label{XXI.7.5.2}
It follows from 7.2.2 and 6.8.4 that $q(\alpha)$ depends only on $\operatorname{long}(\alpha)$. One then easily sees that if $q(\alpha)$ is not constant, $q(\alpha)\operatorname{long}(\alpha)$ is constant, which shows that then $p=2$ or $3$. A glance at the diagrams of the preceding section shows that there are four possible cases (the same letter denotes a Dynkin diagram and the corresponding reduced simply connected root datum):
\[
\begin{array}{lll}
p=2,&B_n\xrightarrow{f_1}C_n,\quad C_n\xrightarrow{f_2}B_n&\text{(with }C_2=B_2\text{).}\\
p=2,&F_4\xrightarrow{g}F_4.\\
p=3,&G_2\xrightarrow{h}G_2.
\end{array}
\]
The reader will note that $f_1\circ f_2$, $f_2\circ f_1$, $g\circ g$ and $h\circ h$ are $p$-morphisms of the form described in 6.8.6.
\end{remark}

\numpara{7.5.3} One sees at once from the preceding description that if $\mathcal R$ and $\mathcal R'$ are two reduced root data of semisimple rank $\leq2$, and one has a $p$-morphism from $\mathcal R'$ to $\mathcal R$, then $\mathcal R$ and $\mathcal R'$ are of the same type. More precisely, one has the following table.

\emph{Notations:} Let $f:\mathcal R'\to\mathcal R$ be a $p$-morphism. Denote by $q$ (resp. $q_1$) any positive power of $p$. For systems of rank~2, use the notations of section~4 (denote the simple roots by $\alpha,\beta$, with $\ell(\alpha)\leq\ell(\beta)$).
\[
\begin{array}{|c|c|c|c|}\hline
\textbf{Type}&p&\text{values of }f&\text{values of }{}^tf\\\hline
\text{Trivial}&\text{arbitrary}&-&-\\\hline
A_1&\text{arbitrary}&f(\alpha')=q\alpha&{}^tf(\alpha^*)=q\alpha'^*\\\hline
A_1\times A_1&\text{arbitrary}&f(\alpha')=q\alpha&{}^tf(\alpha^*)=q\alpha'^*\\
&&f(\beta')=q_1\beta&{}^tf(\beta^*)=q_1\beta'^*\\\hline
A_2,B_2,G_2&\text{arbitrary}&f(\alpha')=q\alpha&{}^tf(\alpha^*)=q\alpha'^*\\
&&f(\beta')=q\beta&{}^tf(\beta^*)=q\beta'^*\\\hline
B_2&p=2&f(\alpha')=q\beta&{}^tf(\alpha^*)=2q\beta'^*\\
&&f(\beta')=2q\alpha&{}^tf(\beta^*)=q\alpha'^*\\\hline
G_2&p=3&f(\alpha')=q\beta&{}^tf(\alpha^*)=3q\beta'^*\\
&&f(\beta')=3q\alpha&{}^tf(\beta^*)=q\alpha'^*\\\hline
\end{array}
\]

\begin{thebibliography}{BLie}
\bibitem[BLie]{XXI.BLie} N.~Bourbaki, \textit{Groupes et alg\`ebres de Lie}, Ch.~IV--VI, Hermann, 1968.
\bibitem[De80]{XXI.De80} M.~Demazure, \textit{A,B,C,D,E,F, etc.}, pp.~221--227 in: S\'eminaire sur les singularit\'es des surfaces (Palaiseau, 1976--1977), eds.~M.~Demazure, H.~C.~Pinkham, B.~Teissier, Lect.~Notes Math.~\textbf{777}, Springer-Verlag, 1980.
\end{thebibliography}
